6.0 PORUSZANIE SIĘ PO LABIRYNCIE

Ogólna zasada:

Gra ,,Labirynt śmierci'' nie posiada planszy w ścisłym tego słowa znaczeniu, gdyż gracze tworzą ją w trakcie gry, konstruując stopniowo labirynt poprzez układanie pasujących do siebie żetonów, oznaczających korytarze i komnaty.
Żeton reprezentujący oddział musi zawsze być umieszczony na jednym z segmentów, oznaczając miejsce pobytu oddziału.
Poruszanie się po terenie labiryntu następuje zawsze z jednego segmentu do segmentu sąsiedniego, tzn. niedopuszczalne jest przejście od razu dwu lub więcej segmentów.

Sposób postępowania:

Gdy oddział zdecyduje się opuścić segment jeden z graczy losowo wybiera z pojemnika następny segment i umieszcza go tak, by przylegał do segmentu opuszczanego.
Należy przy tym pamiętać, by korytarze lub przejścia na obu segmentach łączyły się ze sobą.
Jeżeli istnieje kilka sposobów ułożenia nowego segmentu i każdy z nich spełnia warunek sformułowany w poprzednim zdaniu gracze mają całkowitą swobodę w wyborze jednego z tych sposobów.
Jeżeli oddział porusza się drogą, która już raz przeszedł nie istnieje konieczność ciągnięcia nowych segmentów.

Zasady szczegółowe:

[6.1] Nowy segment musi zostać ułożony w ten sposób w stosunku do wszystkich sąsiednich, by korytarze łączyły się z korytarzami, przejście z przejściami, a ślepe ściany ze ślepymi ścianami.
Jeżeli ułożenie nowego żetonu w ten sposób jest niemożliwe jest on wrzucany z powrotem do pojemnika i losowanie jest powtarzane, nawet kilkakrotnie, tzn. do chwili znalezienia pierwszego pasującego żetonu.

[6.2] Jeżeli istnieje więcej niż jeden sposób ułożenia nowego segmentu gracze decydują w jaki sposób ten segment zostanie ułożony.

[6.3] Oddział zawsze porusza się tylko o jeden segment w czasie jednego etapu gry.

[6.4] Labirynt nie może skończyć się ślepym zaułkiem, dopóki nie zostaną odnalezione Czarne Wrota.
Jeżeli wszystkie wyjścia z aktualnego segmentu wiodą do ślepych zaułków ostatni wybrany żeton jest odkładany z powrotem do pojemnika, zaś ciągnięty jest inny – taki, który nie ,,zamyka'' labiryntu.

[6.5] Labirynt może mieć do trzech poziomów w głąb.

W niektórych komnatach można znaleźć schody.
Każde schody łączą ze sobą wszystkie trzy poziomy (patrz 13.7).
Im głębiej oddział zajdzie, tym spotykane potwory będą groźniejsze, ale również bogatsze, a ich pokonanie będzie dawało więcej punktów doświadczenia.
Patrz 6.7.

[6.6] Przejście z poziomu na inny poziom jest traktowane jak przejście z jednego segmentu do drugiego, a więc nie można przejść od razu, w jednym etapie, dwóch poziomów.
Segment na nowym poziomie jest traktowany tak jak komnata.
Po zejściu oddziału na niższy poziom konstrukcja tego poziomu jest rozpoczynana na nowo od ułożenia żetonu, zawierającego schody.
Uwaga: może się zdarzyć, że na niższym poziomie zostaną znalezione inne schody, niż te, którymi zszedł oddział.
Mogą one zostać włączone w skład labiryntu tylko wtedy, jeżeli istnieje możliwość umieszczenia schodów w odpowiadającym im miejscu również na dwóch pozostałych poziomach, tzn. wtedy, gdy na jednym lub obu pozostałych poziomach w odpowiadającym nowo odnalezionemu schodom miejscu nie został już umieszczony inny żeton.

[6.7] Tabela ,,Poziomy labiryntu'' – patrz plansza z tabelami.

[6.8] Tabela ,,Poziomy labiryntu'' zmienia cechy charakterystyczne potworów, posiadane przez nie skarby oraz zdobywaną za ich zabicie liczbę punktów doświadczenia w zależności od poziomu.
Należy ją odczytywać w następujący sposób:

Siła – dodawana do normalnej siły potwora, np. potwór, którego normalna siła określana jest przez 2+2 (2K6+2), spotkany na poziomie 2+2, będzie mieć Siłę określoną przez 2+4.

Zręczność – dodawana do normalnej zręczności – patrz wyżej.

Niechęć do negocjacji – dodawana do normalnego współczynnika NN – patrz wyżej.

Liczba potworów – należy pomnożyć liczbę napotkanych potworów przez odpowiadającą poziomowi cyfrę, np. jeżeli z tabeli 6.8 wynika, że oddział natnął się na 2 Chimery, a dzieje się to na trzecim poziomie to znaczy, że faktycznie Chimer jest 4 (2 x 2).

Typ skarbu – poza tym wpływa na przesunięcie typu skarbu w dół tabeli 14.9 ,,Skarby'' o jeden stopień dla drugiego poziomu oraz o dwa stopnie dla poziomu trzeciego np. potwór, który normalnie posiadałby skarb typu C, na trzecim poziomie będzie posiadał skarb typu E.

Doświadczenie – liczba punktów doświadczenia, zdobywana przez postać jest mnozona przez cyfrę, odpowiadającą poziomowi.

[6.9] Jeżeli z jakichkolwiek przyczyn potwór, na który oddział się natknął nie został zabity, na który się natknął nie zostało zabity, na który nastąpiło spotkanie umieszcza się odpowiadający potworowi żeton, aby oznaczyć jego pozycję.
Jeżeli oddział jeszcze wejdzie do tego samego segmentu to spotka tego samego potwora i tylko jego, przy czym powinien on być traktowany jako ,,nowy'', tzn. ewentualne poprzednie rezultaty negocjacji czy wykupu są ignorowane.